\subsection{The value at a point and the nearby limit}

After the sequential theory is established, we can separate two questions that concern a function near a point.

\begin{corebox}[Two different questions]
The value
\[
f(a)
\]
asks what the function assigns at the input \(a\). The limit
\[
\lim_{x\to a}f(x)
\]
asks what values are approached along inputs near \(a\), with \(x\ne a\).
\end{corebox}

\begin{definitionbox}[Function limit]
We write
\[
\lim_{x\to a}f(x)=L
\]
when every admissible sequence \((x_n)\) satisfying
\[
x_n\ne a,
\qquad
x_n\to a
\]
produces an output sequence satisfying
\[
f(x_n)\to L.
\]
By Heine's theorem, this sequential statement is equivalent to the usual local definition of a function limit.
\end{definitionbox}

\begin{examplebox}[A removable hole]
For
\[
f(x)=\frac{x^2-1}{x-1},
\qquad x\ne1,
\]
we have
\[
f(x)=x+1
\]
for every admissible \(x\). Therefore, for every sequence \(x_n\to1\) with \(x_n\ne1\),
\[
f(x_n)=x_n+1\to2.
\]
Hence
\[
\lim_{x\to1}f(x)=2,
\]
even though \(f(1)\) is not defined.
\end{examplebox}

\begin{interpretationbox}[The point itself is excluded from the limit test]
Changing, deleting, or redefining the single value \(f(a)\) does not change \(\lim_{x\to a}f(x)\). The limit is determined by values at nearby inputs.
\end{interpretationbox}